\section{Latihan}
\label{sec:latihan-bab1}

\begin{enumerate}
    \item \textbf{Orientasi dokumen.} Buka Daftar Isi buku ini.
          Tuliskan judul bab yang membahas masing-masing topik berikut
          beserta nomor babnya: (a) register \eng{flag},
          (b) mode pengalamatan, (c) prosedur dan stack,
          (d) \eng{disassembly}.

    \item \textbf{Pemetaan capaian.} Jelaskan dengan kalimat Anda
          sendiri perbedaan antara CPL, CPMK, dan Sub-CPMK. Berikan
          satu contoh masing-masing yang diambil dari buku ini.

    \item \textbf{Perancangan mundur.} Pada
          Gambar~\ref{fig:backward-design}, materi berada pada urutan
          terakhir. Mengapa demikian? Apa akibatnya bagi materi yang
          tidak mendukung indikator mana pun?

    \item \textbf{Kesepadanan pertemuan.} Anda sedang berada pada
          minggu ke-6 perkuliahan. Bab mana yang seharusnya Anda baca?
          Sub-CPMK apa yang sedang dibangun, dan CPMK apa yang
          didukungnya?

    \item \textbf{Analisis beban.} Perhatikan
          Tabel~\ref{tab:jalur-capaian}. CPMK mana yang mendapat
          bobot terbesar? Menurut Anda, mengapa bobot itu diberikan
          pada CPMK tersebut?

    \item \textbf{Penilaian kinerja.} Sebutkan tiga perbedaan antara
          cara belajar yang cocok untuk asesmen berbasis kinerja dan
          cara belajar yang cocok untuk asesmen berbasis hafalan.
          Kaitkan jawaban Anda dengan komposisi penilaian mata
          kuliah ini.

    \item \textbf{Penyusunan target.} Tuliskan satu target belajar
          untuk diri Anda sendiri dengan syarat: berkaitan dengan
          salah satu Sub-CPMK, dapat diukur, dan memiliki tenggat
          minggu yang jelas.

    \item \textbf{Perencanaan jurnal.} Rancang satu entri jurnal
          praktikum yang lengkap untuk sesi hari ini sesuai
          Tabel~\ref{tab:format-jurnal}. Isi seluruh kolom dengan
          jujur, termasuk kendala yang Anda alami saat pertama kali
          membuka buku ini.
\end{enumerate}

\begin{refleksibox}
Jawab pertanyaan berikut pada jurnal praktikum Anda.
\begin{enumerate}
    \item Apa harapan Anda terhadap mata kuliah ini sebelum membaca
          bab ini? Apakah bab ini mengubah harapan tersebut?
    \item Bagian mana dari kerangka OBE pada
          Gambar~\ref{fig:backward-design} yang paling meyakinkan
          bagi Anda, dan bagian mana yang masih membingungkan?
    \item Kebiasaan belajar apa yang paling perlu Anda ubah agar
          komponen praktikum yang berbobot \dec{70}\% dapat dicapai
          dengan baik?
\end{enumerate}
\end{refleksibox}
